% =====================================================================
%  ANNEXE G — GLOSSAIRE, TABLE DE REPÉRAGE ET MOT FINAL
% =====================================================================
\chapter{Glossaire, repérage et mot final}

\section{Glossaire des termes grammaticaux}
Tous les termes techniques employés dans ce livre, avec leur équivalent anglais, une
explication en une ligne et un exemple. Classés par ordre alphabétique français.

\begin{tableau}[Glossaire --- 54 termes]
\footnotesize
\begin{longtable}{@{}p{3.6cm}p{3.2cm}p{7.4cm}@{}}
\entete{Terme français} & \entete{Terme anglais} & \entete{Explication et exemple} \tabularnewline
\endfirsthead
\entete{Français} & \entete{Anglais} & \entete{Explication} \tabularnewline
\endhead
\hline
adjectif & \en{adjective} & Mot qui qualifie un nom, placé \textbf{avant} en anglais. \en{a \textbf{beautiful} house}. \tabularnewline
adverbe & \en{adverb} & Modifie un verbe ou un adjectif. \en{She sings \textbf{well}.} \tabularnewline
auxiliaire & \en{auxiliary verb} & Verbe d'aide : \en{be, have, do} et les modaux. \en{She \textbf{has} left.} \tabularnewline
base verbale & \en{base form / bare infinitive} & Forme nue du verbe, sans \en{to} ni \en{-s}. \en{to go \ra go}. \tabularnewline
complément & \en{complement} & Élément qui complète le verbe ou le nom. \en{She is \textbf{a doctor}.} \tabularnewline
comparatif & \en{comparative} & Forme de comparaison entre deux éléments. \en{taller than, more interesting than}. \tabularnewline
compte rendu & \en{minutes / report} & Texte qui résume une réunion ou un événement. \tabularnewline
conditionnel & \en{conditional} & Structure en \en{if} exprimant une hypothèse. \en{If I had time, I'd come.} \tabularnewline
conjonction & \en{conjunction / linking word} & Mot qui relie deux idées. \en{and, but, because, although}. \tabularnewline
contractée (forme) & \en{contraction} & Forme raccourcie de l'oral. \en{I'm, don't, she's, won't, I'd}. \tabularnewline
coordination & \en{coordination} & Jonction de deux éléments de même niveau. \en{tea \textbf{and} coffee}. \tabularnewline
dénombrable & \en{countable} & Nom qui se compte, avec un pluriel. \en{one book, two books}. \tabularnewline
discours rapporté & \en{reported / indirect speech} & Façon de rapporter les paroles d'autrui. \en{She said she was tired.} \tabularnewline
ellipse & \en{ellipsis} & Omission d'un élément sous-entendu. \en{``Who wants tea?'' --- ``Me!''} \tabularnewline
gérondif & \en{gerund} & Forme en \en{-ing} employée comme nom. \en{\textbf{Swimming} is fun. I enjoy \textbf{reading}.} \tabularnewline
impératif & \en{imperative} & Ordre, conseil, invitation. \en{\textbf{Sit} down. \textbf{Don't} worry.} \tabularnewline
indénombrable & \en{uncountable} & Nom sans pluriel, mesuré et non compté. \en{water, information, advice, money}. \tabularnewline
infinitif & \en{infinitive} & \en{to} + base verbale. \en{I want \textbf{to leave}.} \tabularnewline
interrogatif (mot) & \en{Wh-word} & \en{who, what, where, when, why, how, which, whose}. \tabularnewline
intransitif & \en{intransitive} & Verbe sans complément d'objet. \en{He \textbf{arrived}. She \textbf{slept}.} \tabularnewline
marqueur temporel & \en{time marker / signal word} & Mot qui indique le temps à employer. \en{yesterday, since, every day}. \tabularnewline
modal & \en{modal verb} & Auxiliaire de nuance : \en{can, could, may, might, must, should, will, would}. \tabularnewline
nom & \en{noun} & \en{house, freedom, Senegal}. \tabularnewline
négation & \en{negation} & \en{not}, ou les auxiliaires négatifs. \en{I do\textbf{n't} know.} \tabularnewline
objet (complément d') & \en{(direct) object} & Ce qui subit l'action. \en{I read \textbf{a book}.} \tabularnewline
participe passé & \en{past participle} & Troisième forme du verbe. \en{gone, written, eaten, played}. \tabularnewline
participe présent & \en{present participle} & Forme en \en{-ing} des temps continus. \en{I am \textbf{working}.} \tabularnewline
passé simple & \en{past simple / preterite} & Temps du récit achevé. \en{I \textbf{went} home.} \tabularnewline
passif & \en{passive voice} & Sujet qui subit l'action. \en{The bridge \textbf{was built} in 1990.} \tabularnewline
phrase complexe & \en{complex sentence} & Une principale + une ou plusieurs subordonnées. \tabularnewline
phrasal verb & \en{phrasal verb} & Verbe + particule au sens propre. \en{give up, look after, turn off}. \tabularnewline
préposition & \en{preposition} & Mot de relation. \en{in, on, at, for, with, about}. \tabularnewline
présent continu & \en{present continuous} & \en{be + -ing}, action en cours. \en{I \textbf{am reading}.} \tabularnewline
présent parfait & \en{present perfect} & \en{have + participe passé}, passé lié au présent. \en{I \textbf{have seen} it.} \tabularnewline
pronom & \en{pronoun} & Remplace un nom. \en{I, you, he, it, we, they, mine, this}. \tabularnewline
pronom relatif & \en{relative pronoun} & \en{who, which, that, whose, where}. \en{The man \textbf{who} called.} \tabularnewline
proposition & \en{clause} & Groupe sujet + verbe. \en{[I think] [that she is right]}. \tabularnewline
quantifieur & \en{quantifier} & \en{some, any, much, many, a lot of, few, several}. \tabularnewline
question tag & \en{question tag} & Mini-question en fin de phrase. \en{It's cold, \textbf{isn't it}?} \tabularnewline
réfléchi (pronom) & \en{reflexive pronoun} & \en{myself, yourself, himself, ourselves, themselves}. \tabularnewline
registre & \en{register} & Niveau de langue : soutenu, courant, familier. \tabularnewline
relative (proposition) & \en{relative clause} & Subordonnée qui précise un nom. \en{the book \textbf{that I read}}. \tabularnewline
schwa & \en{schwa} & Son neutre très court, le plus fréquent de l'anglais. \en{about \pron{eubaout}}. \tabularnewline
subjonctif & \en{subjunctive} & Rare en anglais, presque invisible. \en{I suggest that he \textbf{be} informed.} \tabularnewline
subordonnée & \en{subordinate clause} & Proposition dépendante. \en{\textbf{When I arrived}, she left.} \tabularnewline
superlatif & \en{superlative} & Le plus haut degré. \en{the tallest, the most interesting}. \tabularnewline
sujet & \en{subject} & Celui qui fait l'action --- \textbf{toujours exprimé} en anglais. \en{\textbf{It} is raining.} \tabularnewline
suffixe & \en{suffix} & Terminaison qui change la catégorie. \en{teach \ra teach\textbf{er}, kind \ra kind\textbf{ness}}. \tabularnewline
syllabe tonique & \en{stressed syllable} & La syllabe prononcée plus fort. \en{ba\textbf{NA}na}. \tabularnewline
temps & \en{tense} & Forme verbale qui situe l'action dans le temps. \tabularnewline
transitif & \en{transitive} & Verbe qui exige un complément d'objet. \en{I \textbf{bought} a car.} \tabularnewline
verbe d'état & \en{stative verb} & Verbe de sensation, opinion, possession --- rarement au \en{-ing}. \en{know, like, want, have, believe}. \tabularnewline
verbe irrégulier & \en{irregular verb} & Verbe aux trois formes non standard. \en{go/went/gone}. \tabularnewline
voix active & \en{active voice} & Sujet qui agit. \en{The dog bit the man.} \tabularnewline
\end{longtable}
\end{tableau}

\section{Table de repérage : où trouver ce que je cherche ?}
\begin{tableau}[Vos questions, nos réponses]
\footnotesize
\begin{longtable}{@{}p{8.2cm}p{5.9cm}@{}}
\entete{Je veux\ldots{}} & \entete{Voir\ldots{}} \tabularnewline
\endfirsthead
\entete{Je veux\ldots{}} & \entete{Voir\ldots{}} \tabularnewline
\endhead
\hline
comprendre comment organiser mon année & chapitre 2 (plan des 12 mois) \tabularnewline
apprendre à travailler efficacement & chapitre 3 (huit techniques) \tabularnewline
trouver des ressources gratuites & chapitre 4 et annexe E \tabularnewline
prononcer un son correctement & chapitres 5 et 6 \tabularnewline
corriger mon accent français & chapitre 7 \tabularnewline
comprendre le \en{present perfect} & chapitre 9 \tabularnewline
maîtriser les articles \en{a/an/the} & chapitre 8 \tabularnewline
poser des questions & chapitre 8 et fiche 3 \tabularnewline
utiliser les modaux & chapitre 10 et fiche 2 \tabularnewline
construire le passif & chapitre 10 et fiche 9 \tabularnewline
employer le discours rapporté & chapitre 10 et fiche 7 \tabularnewline
comprendre les conditionnels & chapitre 11 et fiche 8 \tabularnewline
apprendre du vocabulaire par thème & chapitre 12 \tabularnewline
conjuguer un verbe irrégulier & annexe A \tabularnewline
comprendre un \emph{phrasal verb} & chapitre 13 et annexe D (particules) \tabularnewline
éviter un faux ami & chapitre 13 et fiche 13 \tabularnewline
apprendre les idiomes & chapitre 13 \tabularnewline
dialoguer dans une situation précise & chapitre 14 \tabularnewline
avoir une phrase toute faite & annexe B \tabularnewline
parler sans traduire dans ma tête & chapitre 15 \tabularnewline
me débloquer quand je cherche un mot & chapitre 15 (stratégies de contournement) \tabularnewline
améliorer ma compréhension orale & chapitre 16 \tabularnewline
trouver un texte d'écoute avec corrigé & chapitre 16 (dix textes) \tabularnewline
lire un texte adapté à mon niveau & chapitre 17, partie A \tabularnewline
écrire un e-mail correct & chapitre 17, partie D \tabularnewline
tenir un journal en anglais & chapitre 17, partie E \tabularnewline
connaître l'orthographe britannique ou américaine & chapitre 17, partie B \tabularnewline
corriger mes erreurs typiques & chapitre 18 \tabularnewline
choisir la bonne préposition & annexe D \tabularnewline
varier mes connecteurs & annexe D \tabularnewline
évaluer mon niveau & chapitre 19 \tabularnewline
préparer un examen & chapitre 19 \tabularnewline
réviser toute la grammaire en 14 pages & annexe C \tabularnewline
suivre ma progression au quotidien & annexe F \tabularnewline
\end{longtable}
\end{tableau}

\section{Le mot final}

Vous êtes arrivé au bout de ce livre --- ce qui, statistiquement, est déjà un exploit.
La plupart des gens qui achètent une méthode la feuillettent trois jours. Vous, vous l'avez
lue jusqu'à la dernière annexe.

Mais vous savez maintenant une chose essentielle : \textbf{ce livre n'était pas le but}.
Il était une carte. La carte n'est pas le territoire, et connaître le chemin par c\oe ur ne
fait pas avancer d'un mètre. Ce qui vous fera parler anglais, ce ne sont pas les 3\,000 mots
de ces pages : ce sont les 300 heures pendant lesquelles votre bouche les prononcera, les
50 conversations où vous aurez eu peur et où vous serez resté, les 200 enregistrements que
vous aurez réécoutés en grimaçant, les 400 jours où vous vous serez assis à la même heure
alors que vous n'en aviez aucune envie.

Vous allez connaître trois moments difficiles. Prévoyons-les.

\textbf{Le premier, vers la sixième semaine.} L'enthousiasme retombe, la nouveauté s'use,
et vous aurez l'impression de ne rien apprendre. C'est exactement le moment où les autres
abandonnent. La seule chose à faire : réduire, mais ne pas arrêter. Vingt minutes par jour
plutôt que zéro.

\textbf{Le deuxième, vers le cinquième mois.} Vous comprendrez presque tout, mais vous
parlerez mal, et cet écart entre votre oreille et votre bouche est douloureux. C'est bon
signe : cela signifie que votre compréhension a progressé plus vite que votre production,
ce qui est l'ordre naturel des choses. La parole rattrapera. Continuez les drills du
chapitre 15 --- ils sont faits précisément pour ce passage.

\textbf{Le troisième, vers le neuvième mois.} Vous aurez une conversation ratée. Un mot
manquera au mauvais moment, vous rougirez, et vous vous direz que tout cela ne sert à rien.
Le lendemain, relisez votre journal de progression (annexe F). Regardez la ligne de la
semaine 1. Puis écoutez votre enregistrement du jour 1. Vous ne reconnaîtrez pas cette
personne --- et c'est précisément la preuve que vous êtes devenu quelqu'un d'autre.

Une dernière chose, la plus importante de toutes. Vous n'apprenez pas l'anglais pour passer
un examen, ni pour impressionner quelqu'un, ni pour cocher une case sur un CV. Vous
l'apprenez pour \textbf{rencontrer des gens}. Chaque conversation réelle --- même ratée,
même boiteuse, même avec vingt silences --- vaut plus que dix heures d'étude solitaire.
Allez vers les gens dès le premier mois. Parlez mal, tôt, souvent. C'est le secret que tous
les polyglottes finissent par avouer : ils n'ont jamais été bons avant d'avoir été mauvais
longtemps, en public.

Dans douze mois, quelqu'un vous demandera : « Mais vous parlez anglais ? » Et vous répondrez
--- peut-être en haussant les épaules, parce que vous n'y penserez même plus :
\en{``Yes, more or less. I've been practising for a year.''}

Ce jour-là, pensez à revenir écrire une ligne dans votre journal.

\begin{center}
\vspace{0.6cm}
\textit{\en{``The best time to plant a tree was twenty years ago.\\ The second best time is now.''}}\\[6pt]
\footnotesize\selectlanguage{french}\textit{Le meilleur moment pour planter un arbre,
c'était il y a vingt ans. Le deuxième meilleur moment, c'est maintenant.}
\vspace{0.6cm}
\end{center}

\begin{mission}[Votre première action, dans les dix prochaines minutes]
\textbf{1.} Ouvrez votre téléphone et réglez une alarme quotidienne à l'heure que vous avez
choisie dans votre contrat personnel (annexe F).\\
\textbf{2.} Enregistrez-vous pendant trois minutes : présentez-vous en anglais, comme vous
pouvez, maintenant. Gardez ce fichier. Ne le supprimez jamais.\\
\textbf{3.} Envoyez le message de l'annexe E à une personne.\\
\textbf{4.} Commencez le chapitre 5 demain matin.\\
C'est tout. Quatre gestes, dix minutes, et votre année a commencé.
\end{mission}
